\section*{Chapitre \ref{ch:b3:electrode-kinetics} --- Cinétique électrochimique et électroanalyse}

\begin{solution}{exo:b3:electrode-kinetics:1}
Pente cathodique $2.303RT/\alpha F = \qty{118}{mV}$, d'où $\alpha = 59.2/118 = 0.50$~; $j_0 = \qty{2.0e-6}{A.cm^{-2}}$,
l'ordonnée à l'origine en $\eta = 0$.
\end{solution}

\begin{solution}{exo:b3:electrode-kinetics:2}
$i_0 = \qty{2.0e-4}{A}$~; $R_{\mathrm{ct}} = RT/(Fi_0) = 0.02569/\num{2.0e-4} = \qty{128}{\ohm}$.
\end{solution}

\begin{solution}{exo:b3:electrode-kinetics:3}
Pour \ce{MgSO4}, $I = \frac12(4c + 4c) = 4c$. À \qty{1.0}{mmol/L}, $I = \qty{4.0}{mol.m^{-3}}$ et $\kappa^{-1} =
\qty{0.96}{nm} \times \sqrt{100/4} = \qty{4.8}{nm}$~; à \qty{0.10}{mol/L}, $I = \qty{400}{mol.m^{-3}}$, $\kappa^{-1} = \qty{0.48}{nm}$.
\end{solution}

\begin{solution}{exo:b3:electrode-kinetics:4}
$i_{\mathrm c} = C_{\mathrm{dl}}Av = \num{20e-6} \times 0.0707 \times 0.1 = \qty{0.14}{\micro A}$,
moins de 1\,\% du pic. À
\qty{10}{V/s}, il vaut \qty{14}{\micro A}, alors que le pic ne croît
que jusqu'à $19\sqrt{100} = \qty{190}{\micro A}$~:
le courant capacitif croît comme $v$, le courant faradique comme $\sqrt v$.
\end{solution}

\begin{solution}{exo:b3:electrode-kinetics:5}
$D = \pi\bigl(i\sqrt t/nFAc^*\bigr)^2 = \pi(\num{12.2e-6}/(96\,485 \times 0.0707 \times \num{1.00e-6}))^2 =
\qty{1.0e-5}{cm^2.s^{-1}}$.
\end{solution}

\begin{solution}{exo:b3:electrode-kinetics:6}
$FAc^* = \qty{1.364e-2}{C.cm^{-1}}$, $F/RT = \qty{38.92}{V^{-1}}$. À \qty{50}{mV/s}~: $\sqrt{38.92 \times 0.05 \times \num{7.0e-6}} =
\num{3.69e-3}$, $i_{\mathrm p} = 0.4463 \times \num{1.364e-2} \times \num{3.69e-3} = \qty{22.5}{\micro A}$~;
à \qty{500}{mV/s}, $\sqrt{10}$ fois
plus, \qty{71}{\micro A}.
\end{solution}

\begin{solution}{exo:b3:electrode-kinetics:7}
(a) Réversible. (b) \omterm{def:b3:electrode-kinetics:reversibility}{Quasi réversible}~: l'écart croît avec la vitesse de
balayage.
(c) Une étape chimique suit le transfert d'électrons (EC)~: aux
balayages lents, le produit est consommé avant le balayage retour~; un
balayage rapide prend la chimie de vitesse.
\end{solution}

\begin{solution}{exo:b3:electrode-kinetics:8}
$K_{ij}a_j/a_i = \num{2e-4} \times 0.14/\num{4.0e-3} = 0.007$~: l'électrode surestime de 0.7\,\%.
\end{solution}

\begin{solution}{exo:b3:electrode-kinetics:9}
Moindres carrés~: $b = \qty{0.1093}{\micro A.L.\micro g^{-1}}$, $a = \qty{0.466}{\micro A}$, $s = \qty{0.011}{\micro A}$~; $c_0 = a/b =
\qty{4.26}{\micro g.L^{-1}}$, avec $u(c_0) = \frac sb\sqrt{\frac1n + \frac{\bar y^2}{b^2S_{xx}}} = \qty{0.12}{\micro g.L^{-1}}$
($n = 4$, $\bar y = 1.286$, $S_{xx} = 125$).
\end{solution}

\begin{solution}{exo:b3:electrode-kinetics:10}
$Q = 0.0500 \times 1800 = \qty{90.0}{C}$~; $n(\ce{Cu}) = Q/2F = \qty{4.66e-4}{mol}$,
\qty{29.6}{mg}. Le
résultat ne dépend que de la charge, mesurée à partir d'un courant et
d'une durée, et de la constante de Faraday~: aucun étalon n'est
nécessaire, pourvu que l'électrolyse soit complète et que son rendement
faradique soit de 100\,\%.
\end{solution}

\begin{solution}{exo:b3:electrode-kinetics:11}
$\sqrt{DFv/RT} = \sqrt{\num{1.0e-5} \times 38.92 \times 0.1} = \qty{6.24e-3}{cm.s^{-1}}$~: $\Lambda = 1.6$, \omterm{def:b3:electrode-kinetics:reversibility}{quasi réversible}.
À \qty{10}{V/s}, $\Lambda = 0.16$, encore \omterm{def:b3:electrode-kinetics:reversibility}{quasi réversible} mais plus loin
de la limite réversible~; la vague ne serait réversible qu'au-dessous
d'environ \qty{1}{mV/s}.
\end{solution}

\begin{solution}{exo:b3:electrode-kinetics:12}
$s/b = \qty{0.0838}{mM}$ et $b^2S_{xx} = \qty{203.5}{\micro A^2}$. En $\bar y$~: $0.0838\sqrt{1 + 1/7} = \qty{0.090}{mM}$.
À \qty{18.0}{\micro A}~: $0.0838\sqrt{1.143 + 83.0/203.5} = \qty{0.104}{mM}$.
Trois lectures en $\bar y$~:
$0.0838\sqrt{1/3 + 1/7} = \qty{0.058}{mM}$.
\end{solution}

\begin{solution}{pb:b3:electrode-kinetics:1}
\textbf{1.} $\ce{C6H12O6 -> C6H10O6 + 2 H+ + 2 e-}$ et $\ce{[Fe(CN)6]^3- + e- -> [Fe(CN)6]^4-}$~;
bilan
\[
  \ce{C6H12O6 + 2 [Fe(CN)6]^3- -> C6H10O6 + 2 [Fe(CN)6]^4- + 2 H+} .
\]
À l'électrode, $\ce{[Fe(CN)6]^4- -> [Fe(CN)6]^3- + e-}$.
\textbf{2.} Deux.
\textbf{3.} Le dioxygène dissous dans le sang est en faible quantité et
variable, et son produit, le peroxyde d'hydrogène, n'est oxydé qu'à un
potentiel élevé où de nombreuses autres espèces réagissent~; le
médiateur est présent en excès connu.
\textbf{4.} Pour que sa concentration à la surface soit nulle~: le
courant n'est alors limité que par la diffusion et insensible aux
petites variations de potentiel.
\textbf{5.} Le courant décroît comme $t^{-1/2}$~; les lectures doivent
se faire à l'instant utilisé pour l'étalonnage.
\textbf{6.} Cottrell~: $i\sqrt t = 42.0$, 41.4, 41.7, 41.6, 41.8~:
constant à 1\,\% près.
\textbf{7.} $D = \pi(\text{pente}/FAc)^2 = \pi(\num{41.8e-6}/(96\,485 \times 0.0300 \times \num{1.00e-5}))^2 = \qty{6.55e-6}{cm^2.s^{-1}}$.
\textbf{8.} $\sqrt{\pi \times \num{6.55e-6} \times 5} = \qty{0.010}{cm} = \qty{100}{\micro m}$~:
l'appauvrissement atteint le haut de la couche vers \qty{5}{s}~; des
lectures plus tardives tomberaient au-dessous de la loi de Cottrell.
\textbf{9.} La moitié~: \qty{9.35}{\micro A}.
\textbf{10.} Le courant de fond~: oxydation d'interférents et
d'impuretés, charge de la double couche, hexacyanoferrate(II) résiduel
du médiateur lui-même.
\textbf{11.} Oui~: les résidus, au plus \qty{0.095}{\micro A}, se
dispersent sans tendance. Aux fortes concentrations de glucose, le
médiateur, présent en quantité limitée, ou l'\omterm{def:b3:complex-kinetics:enzyme}{enzyme} se sature et la
droite s'infléchit.
\textbf{12.} $(7.10 - 0.356)/0.9189 = \qty{7.34}{mM}$.
\textbf{13.} $0.0838\sqrt{1 + 1/7 + (7.10 - 8.89)^2/203.5} = 0.0838 \times 1.076 = \qty{0.090}{mM}$.
\textbf{14.} Le terme en $1/m$ (une seule lecture)~; des lectures
répétées, ou plusieurs bandelettes, le réduisent.
\textbf{15.} $3 \times 0.077/0.919 = \qty{0.25}{mM}$.
\textbf{16.} $7.34 \times 180.16 = \qty{1322}{mg.L^{-1}} = \qty{132}{mg.dL^{-1}}$.
\textbf{17.} Il ajoute son propre courant d'oxydation~: le lecteur
surestime.
\textbf{18.} Un \omterm{def:b3:electrode-kinetics:cv}{voltampérogramme} de la bandelette sans glucose montre une
vague anodique de l'ascorbate là où l'hexacyanoferrate(II) est oxydé.
\textbf{19.} Moins de substances sont oxydées à un potentiel plus bas.
\textbf{20.} On retranche le courant de l'électrode témoin de la lecture
de l'électrode à \omterm{def:b3:complex-kinetics:enzyme}{enzyme} avant d'appliquer la droite d'étalonnage.
\textbf{21.} $D$ croît de quelques pour cent par kelvin et le courant
comme $\sqrt D$~: les lecteurs mesurent la température et corrigent en
conséquence.
\textbf{22.} La viscosité et la fraction de globules rouges du sang
modifient $D$ et le courant~; des étalons dans une matrice semblable
compensent ces effets.
\textbf{23.} Les variations d'une bandelette à l'autre (aire de
l'électrode, charge en \omterm{def:b3:complex-kinetics:enzyme}{enzyme} et en médiateur), la température et la
composition du sang s'ajoutent à l'incertitude propre de l'étalonnage.
\textbf{24.} \textbf{$c(\text{glucose}) = 7.34 \pm \qty{0.09}{mmol.L^{-1}}$}
(incertitude-type), soit environ \qty{132}{mg.dL^{-1}}.
\end{solution}
